\chapter{Il secondo rivolto: quando e come usarlo}
\label{chap:quarta_sesta}

Abbiamo visto nel Capitolo \ref{chap:triadi} (\ref{subsec:problema_secondo_rivolto}) il problema acustico dell'accordo in secondo rivolto: la quarta rispetto al basso è un intervallo instabile che necessita di un trattamento specifico. In questo capitolo presentiamo i tre casi in cui il secondo rivolto è non solo ammesso, ma fondamentale nella pratica compositiva.

\section{Il 6/4 cadenziale}
\label{sec:cadenziale}

\begin{definizione}[6/4 cadenziale]
    È l'accordo di I6/4 che precede immediatamente la dominante nella formula cadenzale. Le note al soprano e contralto formano una quarta e una sesta rispetto al basso, che è fermo sulla nota della dominante.
\end{definizione}

\begin{regola}
    \begin{itemize}
        \item Le note in quarta e sesta devono essere \textbf{preparate} (tenute o per grado congiunto).
        \item La quarta risolve scendendo alla terza della dominante.
        \item La sesta risolve scendendo alla quinta della dominante.
        \item Metricamente: il 6/4 cade sul tempo \textbf{forte}, il V cade sul tempo successivo.
    \end{itemize}
\end{regola}

\begin{esempio}[Il 6/4 cadenziale in Do maggiore]
    Progressione IV - I6/4 - V - I.
    \musica{
        \partilate
        { \key do \major \time 4/4 }
        { la'2 sol'4 fa' | mi'2 r }
        { fa'2 mi'4 re' | do'2 r }
        { do'2 do'4 si | sol2 r }
        { fa2 sol2 | do2 r }
        \lyricmode { "IV"2 "I6/4 - V" | "I"2 }
        \figuremode { s2 <6 4>4 <5 3> | s2 }
    }

\end{esempio}

\section{Il 6/4 di passaggio}
\label{sec:64_passaggio}

\begin{definizione}[6/4 di passaggio]
    Si forma quando il basso si muove per grado congiunto attraverso la quinta dell'accordo in un passaggio scalare.
\end{definizione}

\begin{regola}
    \begin{itemize}
        \item Il basso deve muoversi per grado congiunto sia in entrata che in uscita.
        \item L'accordo è metricamente \textbf{debole}.
        \item Le voci superiori procedono per grado congiunto o rimangono ferme.
    \end{itemize}
\end{regola}

\begin{esempio}[6/4 di passaggio]
    Passaggio I - V6/4 - I6 su basso scalare Do-Re-Mi: il tenore resta fermo, il soprano procede per grado congiunto.
    \musica{
        \partilate
        { \key do \major \time 3/4 }
        { do''4 si' do'' }
        { mi'4 re' do' }
        { sol4 sol sol }
        { do4 re mi }
        \lyricmode { "I"4 "V6/4" "I6" }
        \figuremode { s4 <6 4>4 <6>4 }
    }
\end{esempio}

\section{Il 6/4 di pedale (o di ornamento)}
\label{sec:64_pedale}

\begin{definizione}[6/4 di pedale]
    Si forma quando le voci superiori si muovono sopra un basso fermo (pedale).
\end{definizione}

\begin{regola}
    Il basso deve essere consonante con l'accordo che precede e con quello che segue. Il 6/4 funge da ornamento transitorio sopra una nota fissa.
\end{regola}

\section{Sintesi degli usi del secondo rivolto}
\label{sec:sintesi_secondo_rivolto}

\begin{center}
\begin{tabular}{|l|l|l|l|}
\hline
Tipo & Posizione metrica & Movimento del basso & Ruolo \\
\hline
Cadenziale & Forte & Fermo sulla dominante & Preparazione della dominante \\
Passaggio & Debole & Grado congiunto & Collegamento melodico \\
Pedale & Indifferente & Nota fissa & Ornamento armonico \\
\hline
\end{tabular}
\end{center}
